\documentclass[11pt]{article}
\usepackage{notes} % place notes.sty in the same folder
\definecolor{AccN}{HTML}{1E7A45}
\fancyhead[L]{\small\color{NavyN}\nouppercase{\leftmark}}\fancyhead[R]{\small\color{GreyN}Instructor Manual --- Lecture 2}
\usepackage{longtable,multicol,needspace}
\setlength{\emergencystretch}{2em}
\newtcolorbox{instr}[1]{nb,colback=AccN!6,colframe=AccN,title={#1}}
\newtcolorbox{prob}{nb,colback=BluN!5,colframe=BluN,title={Problem statement}}
\newtcolorbox{soln}{nb,colback=GrnN!5,colframe=GrnN,title={Model solution and marking points}}
\newtcolorbox{sayit}{nb,colback=white,colframe=NavyN,title={Read aloud to students}}
\newtcolorbox{live}[1]{nb,colback=AccN!4,colframe=AccN,title={#1}}
\begin{document}
\begin{titlepage}\centering\vspace*{2.2cm}
{\color{AccN}\large\bfseries LECTURE 2\par}\vspace{0.6cm}
{\color{NavyN}\Huge\bfseries Supervised Learning Part 1:\\Classification Algorithms for Defence\par}\vspace{0.6cm}
{\color{AccN}\rule{0.7\textwidth}{1.5pt}}\par\vspace{0.6cm}
{\LARGE Instructor Manual\par}\vspace{0.4cm}
{\large\begin{minipage}{0.85\textwidth}\centering Logistic Regression \textbullet\ kNN \textbullet\ Decision Trees \textbullet\ Random Forest \textbullet\ SVM \textbullet\ Naive Bayes\\Case study: 5-class IADS threat classification, N = 8,500 samples, 11 radar + EW features\end{minipage}\par}\vspace{2cm}
{\large AK DHAMIJA, Researcher \& Instructor, AI \& Machine Learning\par}\vspace{0.2cm}
{\normalsize \href{mailto:dhamija.ak@gmail.com}{dhamija.ak@gmail.com} \quad\textbullet\quad 9811453163\par}\vspace{0.3cm}
{Unit 4: Simulation Strategies in Warfare --- AI \& Machine Learning in Warfare\par}\vfill
\begin{tcolorbox}[colback=NavyN!4,colframe=NavyN,boxrule=0.6pt,arc=2pt,width=0.9\textwidth]\small
\textbf{This manual contains:} the 45-minute session plan with timings; a slide-by-slide teaching guide for all 29 slides; key messages, facilitation tips, anticipated questions with model answers, misconceptions and in-class activities; live demonstrations with the interactive studio; consistency notes on the source material; the complete quiz answer key with rationale and scoring rules; and the assignment with marking scheme and full model solutions.
\end{tcolorbox}
\end{titlepage}
\tableofcontents\newpage


\needspace{12\baselineskip}\section{Lecture at a Glance}

\begin{center}\small\renewcommand{\arraystretch}{1.25}\begin{tabularx}{\textwidth}{@{}>{\bfseries}p{3.3cm}>{\raggedright\arraybackslash}X@{}}\toprule

\textcolor{NavyN}{Lecture}&2 --- Supervised Learning Part 1: Classification Algorithms for Defence\\
\textcolor{NavyN}{Duration}&45 minutes: 5 min opening Q\&A, 30 min teaching, 10 min auto-quiz\\
\textcolor{NavyN}{Slides}&29 (Beamer PDF, 16:9), 16 sections\\
\textcolor{NavyN}{Prerequisite}&Lecture~1; students should have attempted its quiz and assignment.\\
\textcolor{NavyN}{Student materials}&\texttt{Lecture2\_Slides\_Beamer.pdf}; \texttt{Lecture2\_Reference\_Text.pdf}; \texttt{Lecture2\_Quiz.html}; \texttt{Lecture2\_Answers\_Assignment.html} (release after the quiz)\\
\textcolor{NavyN}{Interactive studio}&\texttt{Lecture2\_Classifier\_Playground.html} (live demonstrations in Section~\ref{sec:live})\\
\textcolor{NavyN}{Assessment}&Auto-quiz: 15 MCQs, 40~s each, $+1/-0.25$. Assignment: 5 problems, 100 marks, plus a 20-mark bonus problem.\\\bottomrule\end{tabularx}\end{center}

\subsection*{Overview}Lecture~2 is the first deep dive: supervised \emph{classification} for defence, built around one running dataset --- an integrated air-defence (IADS) track classification problem with five classes and about 8,500 labelled tracks. Six algorithm families are taught with their geometry and mathematics (logistic regression, $k$-NN, decision trees, random forests, gradient boosting, SVMs with kernels, and Naive Bayes), then compared on the same data. The second half of the lecture moves the emphasis from algorithms to \emph{evaluation}: class imbalance, ROC/AUC, precision--recall and $F_\beta$ at the operational base rate. The lecture's message is that model choice matters less than honest evaluation on the class that costs lives.

\subsection*{Learning outcomes}By the end of the lecture students should be able to:\begin{enumerate}\item Explain the decision boundary each classifier produces and why.\item Interpret logistic-regression weights as odds ratios.\item Compute Gini impurity and choose a split; explain bagging, out-of-bag error and boosting.\item Derive the SVM margin $2/\|w\|$, explain the role of $C$ and the kernel trick.\item Diagnose class imbalance and choose appropriate metrics (per-class recall, PR curves, $F_2$).\end{enumerate}

\subsection*{Syllabus mapping}\begin{itemize}\item Supervised learning for defence applications --- \textbf{classification of threat types, target categories and tactical situations}\item Example: \textbf{radar signal classification} (IADS track classification); links to sensor fusion in Lecture~3\item Example: \textbf{satellite image interpretation} (EO/SAR detection, IoU and average precision, change detection, analyst triage)\end{itemize}

\needspace{12\baselineskip}\section{Session Plan (45 Minutes)}

\begin{center}\begin{tikzpicture}[x=0.3444cm,y=1cm]
\fill[AmbN!40](0,0)rectangle(5,1);\node[font=\scriptsize\bfseries,align=center]at(2.5,0.5){Q\&A\\recap};
\fill[AccN!25](5.00,0)rectangle(6.00,1);\draw[white,thick](5.00,0)--(5.00,1);
\fill[NavyN!18](6.00,0)rectangle(7.00,1);\draw[white,thick](6.00,0)--(6.00,1);
\fill[AccN!25](7.00,0)rectangle(9.00,1);\draw[white,thick](7.00,0)--(7.00,1);
\node[font=\tiny,rotate=90,anchor=center]at(8.00,0.5){\parbox{0.9cm}{\centering \S3}};
\fill[NavyN!18](9.00,0)rectangle(11.00,1);\draw[white,thick](9.00,0)--(9.00,1);
\node[font=\tiny,rotate=90,anchor=center]at(10.00,0.5){\parbox{0.9cm}{\centering \S4}};
\fill[AccN!25](11.00,0)rectangle(13.00,1);\draw[white,thick](11.00,0)--(11.00,1);
\node[font=\tiny,rotate=90,anchor=center]at(12.00,0.5){\parbox{0.9cm}{\centering \S5}};
\fill[NavyN!18](13.00,0)rectangle(15.00,1);\draw[white,thick](13.00,0)--(13.00,1);
\node[font=\tiny,rotate=90,anchor=center]at(14.00,0.5){\parbox{0.9cm}{\centering \S6}};
\fill[AccN!25](15.00,0)rectangle(17.00,1);\draw[white,thick](15.00,0)--(15.00,1);
\node[font=\tiny,rotate=90,anchor=center]at(16.00,0.5){\parbox{0.9cm}{\centering \S7}};
\fill[NavyN!18](17.00,0)rectangle(18.00,1);\draw[white,thick](17.00,0)--(17.00,1);
\fill[AccN!25](18.00,0)rectangle(23.00,1);\draw[white,thick](18.00,0)--(18.00,1);
\node[font=\tiny,rotate=90,anchor=center]at(20.50,0.5){\parbox{0.9cm}{\centering \S9}};
\fill[NavyN!18](23.00,0)rectangle(24.00,1);\draw[white,thick](23.00,0)--(23.00,1);
\fill[AccN!25](24.00,0)rectangle(26.00,1);\draw[white,thick](24.00,0)--(24.00,1);
\node[font=\tiny,rotate=90,anchor=center]at(25.00,0.5){\parbox{0.9cm}{\centering \S11}};
\fill[NavyN!18](26.00,0)rectangle(27.00,1);\draw[white,thick](26.00,0)--(26.00,1);
\fill[AccN!25](27.00,0)rectangle(29.00,1);\draw[white,thick](27.00,0)--(27.00,1);
\node[font=\tiny,rotate=90,anchor=center]at(28.00,0.5){\parbox{0.9cm}{\centering \S13}};
\fill[NavyN!18](29.00,0)rectangle(32.00,1);\draw[white,thick](29.00,0)--(29.00,1);
\node[font=\tiny,rotate=90,anchor=center]at(30.50,0.5){\parbox{0.9cm}{\centering \S13}};
\fill[AccN!25](32.00,0)rectangle(33.00,1);\draw[white,thick](32.00,0)--(32.00,1);
\fill[NavyN!18](33.00,0)rectangle(35.00,1);\draw[white,thick](33.00,0)--(33.00,1);
\node[font=\tiny,rotate=90,anchor=center]at(34.00,0.5){\parbox{0.9cm}{\centering \S14}};
\fill[RedN!30](35,0)rectangle(45,1);\node[font=\scriptsize\bfseries,align=center]at(40,0.5){Auto-quiz\\15 Q $\times$ 40 s};
\draw[NavyN](0,0)rectangle(45,1);
\foreach \tt in {0,5,10,15,20,25,30,35,40,45}{\draw[NavyN](\tt,0)--(\tt,-0.12);\node[font=\tiny,below]at(\tt,-0.12){\tt};}
\node[font=\scriptsize,below]at(22.5,-0.45){minutes};
\end{tikzpicture}\end{center}

\begin{center}\small\renewcommand{\arraystretch}{1.2}\begin{tabularx}{\textwidth}{@{}l l X r@{}}\toprule\textbf{Clock} & \textbf{Slides} & \textbf{Segment} & \textbf{Min}\\\midrule

0:00--5:00 & --- & \textbf{Opening Q\&A and recap} (see prompts below) & 5\\

5:00--6:00 & 2 & 1. Objectives \& Data (1 slide) & 1\\

6:00--7:00 & 3 & 2. Framework (1 slide) & 1\\

7:00--9:00 & 4--5 & 3. Decision Boundaries (2 slides) & 2\\

9:00--11:00 & 6--7 & 4. Logistic Regression (2 slides) & 2\\

11:00--13:00 & 8--9 & 5. k-Nearest Neighbours (2 slides) & 2\\

13:00--15:00 & 10--11 & 6. Decision Trees (2 slides) & 2\\

15:00--17:00 & 12--13 & 7. Random Forest (2 slides) & 2\\

17:00--18:00 & 14 & 8. Gradient Boosting (1 slide) & 1\\

18:00--23:00 & 15--17 & 9. Support Vector Machines (3 slides) & 5\\

23:00--24:00 & 18 & 10. Naive Bayes (1 slide) & 1\\

24:00--26:00 & 19--20 & 11. Multi-class \& Imbalance (2 slides) & 2\\

26:00--27:00 & 21 & 12. Algorithm Comparison (1 slide) & 1\\

27:00--29:00 & 22--23 & 13. Evaluation Metrics (2 slides) & 2\\

29:00--32:00 & 24--26 & 13a. Satellite \& SAR Imagery (3 slides) & 3\\

32:00--33:00 & 27 & 13b. Across Disciplines (1 slide) & 1\\

33:00--35:00 & 28--29 & 14. Summary (2 slides) & 2\\

35:00--45:00 & --- & \textbf{Auto-quiz} (instructions 1 min, 15 questions $\times$ 40 s = 10 min of answering; collect scores) & 10\\\bottomrule\end{tabularx}\end{center}

Slide 1 is the title slide and is shown while students settle. The teaching segment covers 28 content slides in 30 minutes, an average pace of 0.9 slides per minute. 

Minute allocations are proportional to slide counts; shift time towards the sections students find hardest.

\begin{instr}{Opening Q\&A prompts (first 5 minutes)}\begin{enumerate}\item From Lecture~1: why can a 99\%-accurate detector produce mostly false alarms?\item What is the difference between training, validation and test data?\item Give one example of shortcut learning.\end{enumerate}Take two or three answers per prompt; correct misconceptions briefly and note any topic that needs revisiting.\end{instr}

\needspace{12\baselineskip}\section{Slide-by-Slide Teaching Guide}

Each subsection gives the slide range, the suggested time, what to emphasise, and the slide titles in order. Numbers refer to pages of the Beamer PDF.

\needspace{8\baselineskip}\subsection{1. Objectives \& Data \normalfont\small\color{GreyN}(slide 2, about 1 min)}

\textbf{Teaching note.} Objectives and the IADS dataset: five classes, feature list (PRF, speed, altitude, RCS, EW activity \ldots), class proportions in training versus the real stream.

\begin{center}\small\renewcommand{\arraystretch}{1.1}\begin{tabularx}{\textwidth}{@{}r X@{}}2 & Lecture 2 --- Objectives \& Dataset Overview\\\end{tabularx}\end{center}

\needspace{8\baselineskip}\subsection{2. Framework \normalfont\small\color{GreyN}(slide 3, about 1 min)}

\textbf{Teaching note.} Recap of the supervised framework: data, hypothesis class, loss, optimisation, evaluation.

\begin{center}\small\renewcommand{\arraystretch}{1.1}\begin{tabularx}{\textwidth}{@{}r X@{}}3 & Quick Recap --- Supervised Learning \& the Classification Framework\\\end{tabularx}\end{center}

\needspace{8\baselineskip}\subsection{3. Decision Boundaries \normalfont\small\color{GreyN}(slides 4--5, about 2 min)}

\textbf{Teaching note.} Show the gallery of decision boundaries; then bias--variance and the train/validation/test protocol with stratified cross-validation.

\begin{center}\small\renewcommand{\arraystretch}{1.1}\begin{tabularx}{\textwidth}{@{}r X@{}}4 & Decision Boundaries --- Each Classifier Carves Space Differently\\
5 & Bias--Variance Trade-off and the Evaluation Protocol\\\end{tabularx}\end{center}

\needspace{8\baselineskip}\subsection{4. Logistic Regression \normalfont\small\color{GreyN}(slides 6--7, about 2 min)}

\textbf{Teaching note.} From linear score to sigmoid probability; cross-entropy loss; weights as log-odds. Work the odds ratio $e^{4.07}\approx58.5$ for EW jamming. Show the XOR failure.

\begin{center}\small\renewcommand{\arraystretch}{1.1}\begin{tabularx}{\textwidth}{@{}r X@{}}6 & Logistic Regression --- From Linear Score to Probability\\
7 & Logistic Regression --- Weights, Interpretation \& Limitations\\\end{tabularx}\end{center}

\needspace{8\baselineskip}\subsection{5. k-Nearest Neighbours \normalfont\small\color{GreyN}(slides 8--9, about 2 min)}

\textbf{Teaching note.} $k$-NN: distance, voting, the need for standardisation, and choosing $k$ by cross-validation ($k=1$: 78.3\%, $k=21$: 90.2\%).

\begin{center}\small\renewcommand{\arraystretch}{1.1}\begin{tabularx}{\textwidth}{@{}r X@{}}8 & k-Nearest Neighbours --- ``You Are Judged by the Company You Keep''\\
9 & kNN --- Choosing $k$ by Cross-Validation on the IADS Data\\\end{tabularx}\end{center}

\needspace{8\baselineskip}\subsection{6. Decision Trees \normalfont\small\color{GreyN}(slides 10--11, about 2 min)}

\textbf{Teaching note.} CART: Gini impurity and information gain, recursive partitioning, depth versus accuracy, extracted rules and feature importance. Compute Gini for a five-class uniform node (0.8).

\begin{center}\small\renewcommand{\arraystretch}{1.1}\begin{tabularx}{\textwidth}{@{}r X@{}}10 & Decision Trees --- CART: Recursive Feature-Space Partitioning\\
11 & Decision Tree --- Depth vs Accuracy, Rules and Feature Importance\\\end{tabularx}\end{center}

\needspace{8\baselineskip}\subsection{7. Random Forest \normalfont\small\color{GreyN}(slides 12--13, about 2 min)}

\textbf{Teaching note.} Bagging plus random feature subsets; why decorrelation reduces variance; out-of-bag error (2.7\%) as free validation; RF results (97.5\%).

\begin{center}\small\renewcommand{\arraystretch}{1.1}\begin{tabularx}{\textwidth}{@{}r X@{}}12 & Random Forest --- Strength in Diversity: Many Weak Learners $\to$ One Strong Classifier\\
13 & Random Forest --- Results on the IADS Data\\\end{tabularx}\end{center}

\needspace{8\baselineskip}\subsection{8. Gradient Boosting \normalfont\small\color{GreyN}(slide 14, about 1 min)}

\textbf{Teaching note.} Boosting fits residuals sequentially; learning rate and number of trees as hyper-parameters; XGBoost regularisation.

\begin{center}\small\renewcommand{\arraystretch}{1.1}\begin{tabularx}{\textwidth}{@{}r X@{}}14 & Gradient Boosting \& XGBoost --- Learning from Residuals\\\end{tabularx}\end{center}

\needspace{8\baselineskip}\subsection{9. Support Vector Machines \normalfont\small\color{GreyN}(slides 15--17, about 5 min)}

\textbf{Teaching note.} Hard margin (margin $2/\|w\|$), soft margin and $C$, then the kernel trick and the RBF kernel. Draw the margin by hand.

\begin{center}\small\renewcommand{\arraystretch}{1.1}\begin{tabularx}{\textwidth}{@{}r X@{}}15 & SVM --- Maximum-Margin Classifier: Hard-Margin Formulation\\
16 & SVM --- Soft Margin \& the $C$ Hyperparameter\\
17 & SVM --- The Kernel Trick: Non-linear Classification via Implicit Feature Maps\\\end{tabularx}\end{center}

\needspace{8\baselineskip}\subsection{10. Naive Bayes \normalfont\small\color{GreyN}(slide 18, about 1 min)}

\textbf{Teaching note.} Bayes' rule, the conditional-independence assumption and why it fails on correlated radar features (83.1\%).

\begin{center}\small\renewcommand{\arraystretch}{1.1}\begin{tabularx}{\textwidth}{@{}r X@{}}18 & Na\"ive Bayes --- Prior $\times$ Likelihood = Posterior\\\end{tabularx}\end{center}

\needspace{8\baselineskip}\subsection{11. Multi-class \& Imbalance \normalfont\small\color{GreyN}(slides 19--20, about 2 min)}

\textbf{Teaching note.} One-vs-rest and one-vs-one; then class imbalance: 92.6\% accuracy with 12.3\% bomber recall on the real stream. Fixes: class weights, SMOTE inside the training fold only, threshold moving.

\begin{center}\small\renewcommand{\arraystretch}{1.1}\begin{tabularx}{\textwidth}{@{}r X@{}}19 & Multi-class Classification Strategies --- Extending Binary Classifiers to 5 Classes\\
20 & Class Imbalance --- When Your Dataset Lies to You\\\end{tabularx}\end{center}

\needspace{8\baselineskip}\subsection{12. Algorithm Comparison \normalfont\small\color{GreyN}(slide 21, about 1 min)}

\textbf{Teaching note.} The comparison table: accuracy, per-class recall, latency, interpretability. No single winner; choose by operational constraint.

\begin{center}\small\renewcommand{\arraystretch}{1.1}\begin{tabularx}{\textwidth}{@{}r X@{}}21 & Algorithm Comparison --- All Classifiers on the IADS Dataset\\\end{tabularx}\end{center}

\needspace{8\baselineskip}\subsection{13. Evaluation Metrics \normalfont\small\color{GreyN}(slides 22--23, about 2 min)}

\textbf{Teaching note.} ROC/AUC (RF 0.998 for the bomber class) and why it can flatter rare-class performance; precision--recall and $F_\beta$ with $\beta=2$ for threat detection.

\begin{center}\small\renewcommand{\arraystretch}{1.1}\begin{tabularx}{\textwidth}{@{}r X@{}}22 & Beyond Accuracy --- ROC \& AUC for Defence AI\\
23 & Precision--Recall, Average Precision and the $F_\beta$ Score\\\end{tabularx}\end{center}

\needspace{8\baselineskip}\subsection{13a. Satellite \& SAR Imagery \normalfont\small\color{GreyN}(slides 24--26, about 3 min)}

\textbf{Teaching note.} Satellite and SAR image interpretation (syllabus example): the four tasks (classification, detection, segmentation, change detection), tiling huge scenes, one- versus two-stage detectors and NMS. Work the IoU example (1,225/1,975 = 0.62, a hit) and the average-precision example (AP = 0.68 for 8 detections of 5 vehicles). Close with SAR look-angle shift and the analyst-triage loop: AI ranks, analysts decide.

\begin{center}\small\renewcommand{\arraystretch}{1.1}\begin{tabularx}{\textwidth}{@{}r X@{}}24 & Satellite \& SAR Image Interpretation --- Four Tasks, One CNN Backbone\\
25 & Detection Metrics --- Intersection over Union and Average Precision\\
26 & SAR Is Not a Photograph --- Pitfalls and the Imagery-Triage Concept\\\end{tabularx}\end{center}

\needspace{8\baselineskip}\subsection{13b. Across Disciplines \normalfont\small\color{GreyN}(slide 27, about 1 min)}

\textbf{Teaching note.} Classification in medicine (triage), psychology (signal-detection theory) and aeronautics (fault diagnosis) --- same mathematics, different costs.

\begin{center}\small\renewcommand{\arraystretch}{1.1}\begin{tabularx}{\textwidth}{@{}r X@{}}27 & Classification Across the Disciplines --- From Radar Scopes to Hospital Wards\\\end{tabularx}\end{center}

\needspace{8\baselineskip}\subsection{14. Summary \normalfont\small\color{GreyN}(slides 28--29, about 2 min)}

\textbf{Teaching note.} Summary and look-ahead to regression and sensor fusion.

\begin{center}\small\renewcommand{\arraystretch}{1.1}\begin{tabularx}{\textwidth}{@{}r X@{}}28 & Lecture 2 Summary --- Six Classifiers, One Mission\\
29 & Looking Ahead --- Lecture 3: Regression \& Multi-Sensor Fusion\\\end{tabularx}\end{center}

\needspace{12\baselineskip}\section{Key Messages}Students should leave the room able to state these in their own words:\begin{keypoint}\begin{enumerate}\item Every classifier is a way of carving feature space; the shape of the boundary follows from the model's assumptions.\item Ensembles (random forest, boosting) are strong defaults for tabular defence data.\item Accuracy is the wrong headline on imbalanced data; report per-class recall and precision at the operational base rate.\item Resampling and feature scaling must be fitted on the training fold only, or results leak.\item Choose the model by the operational constraint (latency, explainability, cost of a missed bomber), not by a leaderboard.\end{enumerate}\end{keypoint}

\needspace{12\baselineskip}\section{Facilitation Tips}\begin{itemize}\item Keep the IADS dataset as a thread: every algorithm returns to the same five classes, so students can compare like with like.\item Draw one 2-D toy dataset on the board and sketch each boundary on it as you go: linear, jagged ($k$-NN), axis-aligned boxes (tree), smoothed boxes (forest), curved (RBF SVM).\item The imbalance slide is the turning point; pause there and ask, ``Would you deploy the 92.6\% model?''\item If time is short, compress Naive Bayes and multi-class strategies to one minute each.\end{itemize}

\needspace{12\baselineskip}\section{Anticipated Questions and Model Answers}

\begin{question}{anticipated question}\textbf{Q:} Why not always use the most accurate model?\par\smallskip\textbf{A:} Because accuracy hides the rare class, and operational constraints (latency, explainability, robustness) differ. A 3~ms latency budget or a need for auditable rules may favour a shallower model.\end{question}

\begin{question}{anticipated question}\textbf{Q:} Is SMOTE cheating?\par\smallskip\textbf{A:} Not if it is applied only to the training fold. Applying it before the split leaks synthetic copies of test points into training and inflates results.\end{question}

\begin{question}{anticipated question}\textbf{Q:} Why does a random forest not overfit as trees are added?\par\smallskip\textbf{A:} Averaging more decorrelated trees reduces variance without increasing bias; the error converges. Individual trees overfit, the average does not (much).\end{question}

\begin{question}{anticipated question}\textbf{Q:} How do I pick $C$ and $\gamma$ for an RBF SVM?\par\smallskip\textbf{A:} Grid or random search with stratified cross-validation, on standardised features, scoring with the metric that matters (e.g.\ bomber recall or $F_2$).\end{question}

\needspace{12\baselineskip}\section{Common Misconceptions}\begin{center}\small\renewcommand{\arraystretch}{1.35}\begin{tabularx}{\textwidth}{@{}>{\raggedright\arraybackslash}X >{\raggedright\arraybackslash}X@{}}\toprule\textbf{Misconception} & \textbf{Correction}\\\midrule``AUC near 1 means the model is ready.'' & At a 0.8\% base rate precision can still be poor; check the PR curve.\\
``Feature importance shows causal drivers.'' & Importance is predictive association, often split among correlated features.\\
``More trees in boosting always help.'' & Boosting can overfit; use early stopping.\\
``Naive Bayes is useless because its assumption is false.'' & It is fast and often well-ranked; it suffers when correlated features double-count evidence.\\\bottomrule\end{tabularx}\end{center}

\needspace{12\baselineskip}\section{In-Class Activities and Discussion Prompts}\label{sec:act}\begin{itemize}\item \textbf{Gini at the board:} split a node of 10 tracks (6 bombers, 4 fighters) on ``EW active?'' and compute the impurity decrease.\item \textbf{Deploy or not?} Show the 92.6\%/12.3\% model and ask teams to write one sentence each to the IADS commander recommending action.\end{itemize}

\needspace{12\baselineskip}\section{Run It Live: Studio Demonstrations}\label{sec:live}

Project \texttt{Lecture2\_Classifier\_Playground.html} (it is a single HTML file and runs offline in any browser; every result updates as an input changes). Each demonstration below follows the same rhythm: \textbf{set} the inputs, \textbf{ask} the class to commit to a prediction (a show of hands or a one-word answer), then \textbf{reveal} the result and \textbf{land} the point. The numbers are exactly what the page shows for the settings given; simulations use fixed random seeds, so they repeat every time. The 45-minute plan has little slack: choose the one demonstration marked \emph{first choice} for class, and use the others in a tutorial, a second session, or as homework with the drill cards at the bottom of the page.

\begin{live}{Demo 1: Who sets the threshold? (about 4 min, with \S13 Evaluation Metrics) --- \emph{first choice}}\textbf{Set.} Demo~1 at its defaults: $d'=2.5$, threshold 1.65, 8 threats per 1,000 tracks, a miss costs 20 times a false alarm. The page shows recall 80.2\%, precision 11.6\%, AUC 0.961.\par\smallskip\textbf{Ask.} ``Misses cost twenty times as much as false alarms. Should the threshold go up or down?''\par\smallskip\textbf{Reveal.} Press \textsf{Move to cost-optimal threshold}: it goes \emph{up}, to \textbf{2.00}. Recall falls to 69.1\%, precision rises to 19.7\%, and false alarms per 10,000 tracks fall from 490.8 to 225.7. Now type 80 threats per 1,000 and press the button again: the optimum drops to \textbf{1.05}, with recall 92.6\% and precision 35.4\%.\par\smallskip\textbf{Land the point.} With rare threats, the sheer number of false alarms outweighs even a 20:1 cost ratio. The best threshold depends on base rate and costs, both command inputs; the AUC never changed.\end{live}

\begin{live}{Demo 2: Pick k for k-NN (about 2 min, with \S5 k-Nearest Neighbours)}\textbf{Set.} Demo~2 with $k=1$.\par\smallskip\textbf{Ask.} ``Training accuracy is 100\%. Is $k=1$ the model to deploy?''\par\smallskip\textbf{Reveal.} $k=1$: training 100\%, leave-one-out \textbf{83\%}. $k=5$: 93\% and \textbf{88\%} (the best). $k=41$: 67\% and 67\%.\par\smallskip\textbf{Land the point.} Training accuracy rewards memorising. Choose $k$ on held-out data: too small overfits the noise, too large washes out the boundary.\end{live}

\begin{live}{Demo 3: A near miss in satellite imagery (about 2 min, with \S13a Satellite \& SAR Imagery)}\textbf{Set.} Demo~4 defaults: a 40-pixel vehicle box, prediction offset by 5 pixels; IoU 0.620, a hit at the 0.5 threshold.\par\smallskip\textbf{Ask.} ``Shift the prediction 15 pixels down and right. Still a hit?''\par\smallskip\textbf{Reveal.} Set the predicted box to (25, 25, 65, 65): intersection 625, union 2,575, \textbf{IoU 0.243}, a miss.\par\smallskip\textbf{Land the point.} For small objects a few pixels of error destroys the overlap. At typical ground sample distances that is a few metres, so localisation accuracy matters as much as detection.\end{live}

\needspace{12\baselineskip}\section{Quiz Administration and Answer Key}

\begin{sayit}``You will now take a 15-question quiz. Each question has four options and exactly one correct answer. You have 40 seconds per question; the quiz advances automatically when time runs out. A correct answer scores $+1$, a wrong answer $-0.25$, and an unanswered question 0. The maximum is 15 and the minimum is $-3.75$. Because of the negative marking, guessing blindly among four options gains nothing on average --- answer when you can eliminate options. Your score and a full review appear at the end.''\end{sayit}

\textbf{Delivery.} Open \texttt{Lecture2\_Quiz.html} on each student device (or project it and collect answers on paper). The page enforces the timer and computes the score with negative marking. Release \texttt{Lecture2\_Answers\_Assignment.html} only after all students have finished.

\subsection*{Answer key}\begin{center}\small\renewcommand{\arraystretch}{1.3}\begin{tabular}{|c|c|c|c|c|c|c|c|c|c|c|c|c|c|c|}\hline \textbf{1} & \textbf{2} & \textbf{3} & \textbf{4} & \textbf{5} & \textbf{6} & \textbf{7} & \textbf{8} & \textbf{9} & \textbf{10} & \textbf{11} & \textbf{12} & \textbf{13} & \textbf{14} & \textbf{15}\\\hline B & A & C & A & C & C & D & B & B & D & D & A & A & C & D\\\hline\end{tabular}\end{center}

\textbf{Interpreting scores.} 12 or more: secure command of the lecture. 9--11.75: good, review missed items. Below 9: revisit the reference text before the assignment. If more than half the class misses a question, revisit that point in the next lecture's opening Q\&A.

\subsection*{Questions, options and rationale}

{\small

\noindent\textbf{Q1.} Logistic regression's decision boundary is:\begin{itemize}[leftmargin=2.2em,itemsep=0pt,topsep=1pt]\item[A.] A set of axis-aligned steps\item[\textbf{B.}] \textbf{A hyperplane $w^\top x+b=0$}~{\color{GrnN}$\checkmark$}\item[C.] Always curved\item[D.] A Voronoi tessellation\end{itemize}\noindent{\color{GreyN}\textit{Rationale:}} $\sigma(z)\ge0.5\iff z\ge0$, so the boundary is the hyperplane $w^\top x+b=0$.\par\medskip

\noindent\textbf{Q2.} The EW-jamming weight for the bomber class is $w_9=4.07$. The odds ratio is about:\begin{itemize}[leftmargin=2.2em,itemsep=0pt,topsep=1pt]\item[\textbf{A.}] \textbf{58.5}~{\color{GrnN}$\checkmark$}\item[B.] 4.07\item[C.] 0.98\item[D.] 16.6\end{itemize}\noindent{\color{GreyN}\textit{Rationale:}} Odds ratio $=e^{w}=e^{4.07}\approx58.5$: jamming multiplies the odds of 'bomber' 58-fold.\par\medskip

\noindent\textbf{Q3.} Why can't logistic regression separate low-slow UAVs from civilian helicopters in an XOR-like layout?\begin{itemize}[leftmargin=2.2em,itemsep=0pt,topsep=1pt]\item[A.] Too little data\item[B.] The sigmoid saturates\item[\textbf{C.}] \textbf{A single hyperplane cannot separate interleaved classes}~{\color{GrnN}$\checkmark$}\item[D.] Features are not scaled\end{itemize}\noindent{\color{GreyN}\textit{Rationale:}} Linear models need interaction features or non-linear methods (RF, RBF-SVM) for XOR-type structure.\par\medskip

\noindent\textbf{Q4.} A predicted vehicle box and the true box have an intersection of 1,225 px$^2$ and a union of 1,975 px$^2$. At the usual IoU threshold of 0.5 the detection is:\begin{itemize}[leftmargin=2.2em,itemsep=0pt,topsep=1pt]\item[\textbf{A.}] \textbf{A hit: IoU $\approx$ 0.62}~{\color{GrnN}$\checkmark$}\item[B.] A miss: IoU $\approx$ 0.38\item[C.] A miss, because the boxes are not identical\item[D.] It cannot be judged without the class label\end{itemize}\noindent{\color{GreyN}\textit{Rationale:}} IoU = intersection/union = 1,225/1,975 = 0.62 $\geq$ 0.5, so it counts as a correct detection.\par\medskip

\noindent\textbf{Q5.} Which algorithms require feature standardisation?\begin{itemize}[leftmargin=2.2em,itemsep=0pt,topsep=1pt]\item[A.] Decision trees and random forests\item[B.] Naive Bayes only\item[\textbf{C.}] \textbf{kNN and SVM}~{\color{GrnN}$\checkmark$}\item[D.] None\end{itemize}\noindent{\color{GreyN}\textit{Rationale:}} Distance- and inner-product-based methods are dominated by large-scale features; trees are scale-invariant.\par\medskip

\noindent\textbf{Q6.} A node has five equal classes. Its Gini impurity is:\begin{itemize}[leftmargin=2.2em,itemsep=0pt,topsep=1pt]\item[A.] 0\item[B.] 0.5\item[\textbf{C.}] \textbf{0.8}~{\color{GrnN}$\checkmark$}\item[D.] 1.0\end{itemize}\noindent{\color{GreyN}\textit{Rationale:}} $1-5(0.2^2)=0.8$.\par\medskip

\noindent\textbf{Q7.} Random forests reduce variance mainly by:\begin{itemize}[leftmargin=2.2em,itemsep=0pt,topsep=1pt]\item[A.] Pruning each tree\item[B.] Using a larger learning rate\item[C.] Sequentially fitting residuals\item[\textbf{D.}] \textbf{Averaging many trees made less correlated through bootstrap samples and random feature subsets}~{\color{GrnN}$\checkmark$}\end{itemize}\noindent{\color{GreyN}\textit{Rationale:}} $\text{Var}=\rho\sigma^2+(1-\rho)\sigma^2/B$: averaging removes the second term; feature sub-sampling lowers $\rho$.\par\medskip

\noindent\textbf{Q8.} A SAR vehicle classifier scores 99\% on a benchmark recorded at one depression angle but fails on imagery taken at another angle. The main cause is:\begin{itemize}[leftmargin=2.2em,itemsep=0pt,topsep=1pt]\item[A.] Too few convolutional layers\item[\textbf{B.}] \textbf{Distribution shift: SAR signatures change strongly with look geometry}~{\color{GrnN}$\checkmark$}\item[C.] Wrong labels in the benchmark\item[D.] SAR images carry no useful shape information\end{itemize}\noindent{\color{GreyN}\textit{Rationale:}} SAR returns depend on aspect and depression angle; test on the deployment sensor geometry, not only on the benchmark.\par\medskip

\noindent\textbf{Q9.} Boosting differs from bagging because boosting:\begin{itemize}[leftmargin=2.2em,itemsep=0pt,topsep=1pt]\item[A.] Trains trees independently in parallel\item[\textbf{B.}] \textbf{Trains trees sequentially, each fitting the residuals of the ensemble so far}~{\color{GrnN}$\checkmark$}\item[C.] Never overfits\item[D.] Uses no learning rate\end{itemize}\noindent{\color{GreyN}\textit{Rationale:}} $F_m=F_{m-1}+\eta T_m$, with $T_m$ fitted to $y-F_{m-1}(x)$.\par\medskip

\noindent\textbf{Q10.} In a hard-margin SVM the margin width is:\begin{itemize}[leftmargin=2.2em,itemsep=0pt,topsep=1pt]\item[A.] $\|w\|$\item[B.] $C$\item[C.] $1/b$\item[\textbf{D.}] \textbf{$2/\|w\|$}~{\color{GrnN}$\checkmark$}\end{itemize}\noindent{\color{GreyN}\textit{Rationale:}} Minimising $\tfrac12\|w\|^2$ maximises the margin $2/\|w\|$.\par\medskip

\noindent\textbf{Q11.} Increasing the SVM soft-margin parameter C:\begin{itemize}[leftmargin=2.2em,itemsep=0pt,topsep=1pt]\item[A.] Widens the margin and tolerates more violations\item[B.] Removes support vectors\item[C.] Has no effect with an RBF kernel\item[\textbf{D.}] \textbf{Narrows the margin, penalising violations more (risk of overfitting)}~{\color{GrnN}$\checkmark$}\end{itemize}\noindent{\color{GreyN}\textit{Rationale:}} Large C approaches the hard margin; the IADS optimum was C=100.\par\medskip

\noindent\textbf{Q12.} The kernel trick works because the SVM dual depends on the data only through:\begin{itemize}[leftmargin=2.2em,itemsep=0pt,topsep=1pt]\item[\textbf{A.}] \textbf{Inner products $x_i^\top x_j$}~{\color{GrnN}$\checkmark$}\item[B.] Labels\item[C.] Feature means\item[D.] Distances to the centroid\end{itemize}\noindent{\color{GreyN}\textit{Rationale:}} Replace $x_i^\top x_j$ by $K(x_i,x_j)=\phi(x_i)^\top\phi(x_j)$ without computing $\phi$.\par\medskip

\noindent\textbf{Q13.} Naive Bayes performed worst on IADS (83.1\%) mainly because:\begin{itemize}[leftmargin=2.2em,itemsep=0pt,topsep=1pt]\item[\textbf{A.}] \textbf{Features such as PRF and speed are strongly correlated, violating conditional independence}~{\color{GrnN}$\checkmark$}\item[B.] Too few classes\item[C.] It overfits\item[D.] It needs scaled features\end{itemize}\noindent{\color{GreyN}\textit{Rationale:}} The independence assumption double-counts correlated evidence.\par\medskip

\noindent\textbf{Q14.} On the real IADS stream (0.8\% bombers), a model with 92.6\% accuracy had bomber recall 12.3\%. The best single fix is:\begin{itemize}[leftmargin=2.2em,itemsep=0pt,topsep=1pt]\item[A.] Report accuracy only\item[B.] Delete civilian tracks from the test set\item[\textbf{C.}] \textbf{Class weighting or SMOTE on the training fold, plus per-class recall reporting}~{\color{GrnN}$\checkmark$}\item[D.] Increase k\end{itemize}\noindent{\color{GreyN}\textit{Rationale:}} Weighted loss ($w_{AB}=26.9$) or SMOTE raised bomber recall to about 89\%; always split before oversampling.\par\medskip

\noindent\textbf{Q15.} For a rare, high-cost threat class, the most informative evaluation is:\begin{itemize}[leftmargin=2.2em,itemsep=0pt,topsep=1pt]\item[A.] Overall accuracy\item[B.] ROC-AUC on a balanced lab set only\item[C.] Training accuracy\item[\textbf{D.}] \textbf{Precision-recall / average precision and $F_2$ at the operational base rate}~{\color{GrnN}$\checkmark$}\end{itemize}\noindent{\color{GreyN}\textit{Rationale:}} With rare positives, FPR hides the volume of false alarms; $F_2$ weights recall twice as much as precision.\par\medskip

}

\needspace{12\baselineskip}\section{Assignment: Marking Scheme and Model Solutions}

\textbf{Instructions.} Five problems, 100 marks, about 3 hours, plus a 20-mark bonus problem (Problem 6). Show all working. Reference: Lecture 2 reference text.

\begin{center}\small\renewcommand{\arraystretch}{1.25}\begin{tabularx}{0.85\textwidth}{@{}l X r@{}}\toprule\textbf{\#} & \textbf{Problem} & \textbf{Marks}\\\midrule1 & Metrics on the operational stream & 20\\
2 & Choosing a tree split & 20\\
3 & SVM geometry and slack & 20\\
4 & Bayes with realistic priors & 20\\
5 & Ensembles and imbalance & 20\\
6 & Problem 6 (bonus): Scoring a vehicle detector on satellite imagery & 20\\\midrule & \textbf{Total} & \textbf{100} + 20 bonus\\\bottomrule\end{tabularx}\end{center}

\begin{instr}{General marking guidance}\begin{itemize}
\item Marks for each part are shown in square brackets in the problem statement.
\item Award method marks for correct working even when an arithmetic slip carries through; deduct once for the slip, not again for its consequences.
\item For interpretive parts, full marks need (i) the correct conclusion, (ii) the reason, and (iii) an operational implication or caveat.
\item Answers that state a number without uncertainty, baseline or caveat where one is required lose a mark per part.
\item Suggested deadline: one week after the lecture; return marked work before the next quiz on the same topic.
\end{itemize}\end{instr}

\needspace{8\baselineskip}\subsection{Problem 1: Metrics on the operational stream \normalfont\small\color{GreyN}(20 marks)}

\begin{prob}On 10,000 real tracks containing 80 adversary bombers, a classifier gives TP = 72, FN = 8, FP = 150, TN = 9,770 for the bomber class.

(a) Compute accuracy, precision, recall, $F_1$ and $F_2$. [8]\newline (b) Lowering the threshold gives TP = 76, FP = 400. Recompute precision, recall and $F_2$. [6]\newline (c) Which threshold would you recommend for an IADS, and why might $F_2$ disagree with an operator's view? [6]\end{prob}

\begin{soln}\textbf{(a)} Accuracy $=(72+9770)/10000=98.4\%$. $P=72/222=0.324$; $R=72/80=0.900$; $F_1=2(0.324)(0.9)/(1.224)=0.477$; $F_2=5PR/(4P+R)=1.458/2.196=\mathbf{0.664}$. The 98.4\% accuracy hides that two of every three alarms are false.

\textbf{(b)} $P=76/476=0.160$; $R=76/80=0.950$; $F_2=5(0.160)(0.95)/(0.640+0.95)=0.760/1.590=\mathbf{0.478}$.

\textbf{(c)} By $F_2$ the original threshold is better (0.664 vs 0.478): 4 extra bombers cost 250 extra false alarms. An operator may still prefer the lower threshold if a single missed bomber is catastrophic and a second sensor can cheaply clear false alarms; the right choice depends on explicit costs, workload and the availability of a second-stage check, which $F_\beta$ only approximates.\end{soln}

\needspace{8\baselineskip}\subsection{Problem 2: Choosing a tree split \normalfont\small\color{GreyN}(20 marks)}

\begin{prob}A node holds 100 tracks: 40 AB and 60 FF. Split A gives left (30 AB, 10 FF) and right (10 AB, 50 FF). Split B gives left (40 AB, 30 FF) and right (0 AB, 30 FF).

(a) Compute the parent Gini impurity. [3]\newline (b) Compute the weighted child impurity and Gini reduction for each split. [10]\newline (c) Which split does CART choose? Comment on why split B, despite a pure child, loses. [7]\end{prob}

\begin{soln}\textbf{(a)} $G=1-(0.4^2+0.6^2)=0.48$.

\textbf{(b)} Split A: left $G=1-(0.75^2+0.25^2)=0.375$; right $G=1-((1/6)^2+(5/6)^2)=0.278$; weighted $0.4(0.375)+0.6(0.278)=0.317$; reduction $\mathbf{0.163}$. Split B: left $G=1-((4/7)^2+(3/7)^2)=0.490$; right $G=0$; weighted $0.7(0.490)=0.343$; reduction $\mathbf{0.137}$.

\textbf{(c)} CART chooses A. Split B purifies only 30 FF tracks and leaves a large, very mixed left node (70 tracks, impurity 0.49). Split A makes both children substantially purer. Impurity reduction weights each child by its size, so a small pure child does not compensate for a large impure one.\end{soln}

\needspace{8\baselineskip}\subsection{Problem 3: SVM geometry and slack \normalfont\small\color{GreyN}(20 marks)}

\begin{prob}A linear SVM has $w=(3,4)$, $b=-10$; classes $y\in\{-1,+1\}$.

(a) Compute the margin width. [4]\newline (b) Classify $x_1=(4,2)$ and $x_2=(1,1)$ and compute each point's distance to the boundary. [6]\newline (c) Compute the slack $\xi$ for $(2,1)$ with $y=+1$ and for $(3,1)$ with $y=-1$, and interpret. [6]\newline (d) Explain how increasing $C$ would change the solution. [4]\end{prob}

\begin{soln}\textbf{(a)} $\|w\|=5$; margin $=2/5=\mathbf{0.4}$.

\textbf{(b)} $f(x_1)=12+8-10=10>0\Rightarrow+1$; distance $10/5=2.0$. $f(x_2)=3+4-10=-3<0\Rightarrow-1$; distance $3/5=0.6$.

\textbf{(c)} $(2,1)$, $y=+1$: $f=6+4-10=0$, $y f=0$, $\xi=\max(0,1-0)=\mathbf{1}$ --- exactly on the boundary. $(3,1)$, $y=-1$: $f=9+4-10=3$, $yf=-3$, $\xi=1-(-3)=\mathbf{4}$ --- misclassified ($\xi>1$) and far on the wrong side.

\textbf{(d)} Larger $C$ penalises slack more: the optimiser narrows the margin (larger $\|w\|$) to reduce violations, fitting the training data more closely at greater risk of overfitting; very large $C$ approaches the hard-margin SVM.\end{soln}

\needspace{8\baselineskip}\subsection{Problem 4: Bayes with realistic priors \normalfont\small\color{GreyN}(20 marks)}

\begin{prob}PRF is Gaussian: FF $\mu=1040,\sigma=185$; AB $\mu=2280,\sigma=310$ (Hz). A track has PRF 1,950 Hz.

(a) Compute the two class-conditional densities. [8]\newline (b) Compute $P(\text{AB}\mid x)$ with equal priors and with operational priors $P(\text{AB})=0.01$, $P(\text{FF})=0.99$. [8]\newline (c) A second feature, speed, is strongly correlated with PRF for bombers. Explain what Naive Bayes will do wrong if it is added. [4]\end{prob}

\begin{soln}\textbf{(a)} FF: $z=4.92$, $\phi(z)=0.3989e^{-12.10}=2.21\times10^{-6}$, density $=\phi/\sigma=1.20\times10^{-8}$. AB: $z=-1.065$, $\phi=0.3989e^{-0.567}=0.226$, density $=0.226/310=7.30\times10^{-4}$.

\textbf{(b)} Equal priors: $P(\text{AB}|x)=7.30\times10^{-4}/(7.30\times10^{-4}+1.20\times10^{-8})\approx\mathbf{0.99998}$. Operational priors: numerators $0.01(7.30\times10^{-4})=7.30\times10^{-6}$ and $0.99(1.20\times10^{-8})=1.19\times10^{-8}$, so $P(\text{AB}|x)\approx\mathbf{0.998}$. The evidence is so strong here that even a 1\% prior leaves the posterior high; for weaker evidence the prior would dominate (base-rate effect).

\textbf{(c)} Naive Bayes multiplies the two likelihoods as if independent, effectively counting the same evidence twice. Posteriors become over-confident, and when the features disagree in unusual ways the classifier can be confidently wrong. Remedies: drop one feature, decorrelate (e.g. PCA), or model the joint distribution.\end{soln}

\needspace{8\baselineskip}\subsection{Problem 5: Ensembles and imbalance \normalfont\small\color{GreyN}(20 marks)}

\begin{prob}(a) Trees have prediction variance $\sigma^2=0.04$ and pairwise correlation $\rho=0.3$. Compute the ensemble variance for $B=1,10,500$ and the limit as $B\to\infty$. What would reduce the limit? [8]\newline (b) A training set has 5,000 civilian, 600 friendly, 120 UAV, 40 missile and 20 bomber tracks. Compute the class weights $w_c=N/(C\,N_c)$. [6]\newline (c) A colleague applies SMOTE to the full dataset, then splits into train and test, and reports bomber recall 0.99. Explain the flaw and the correct procedure. [6]\end{prob}

\begin{soln}\textbf{(a)} $\text{Var}=\rho\sigma^2+(1-\rho)\sigma^2/B=0.012+0.028/B$. $B=1$: 0.040; $B=10$: 0.0148; $B=500$: 0.012056; limit 0.012. Only lowering $\rho$ (more randomisation: smaller $m$, sub-sampling) or $\sigma^2$ reduces the floor; more trees cannot.

\textbf{(b)} $N=5780$, $C=5$, $N/C=1156$. Civilian $1156/5000=0.23$; friendly $1.93$; UAV $9.63$; missile $28.9$; bomber $\mathbf{57.8}$.

\textbf{(c)} Synthetic test bombers are interpolations of training bombers, so the test set leaks training information and recall is grossly inflated. Correct: split first (stratified), apply SMOTE (or class weights) only inside the training folds of cross-validation, and evaluate on untouched real test data, ideally at the operational class mix.\end{soln}

\needspace{8\baselineskip}\subsection{Problem 6 (bonus): Scoring a vehicle detector on satellite imagery \normalfont\small\color{GreyN}(20 marks)}

\begin{prob}A scene contains 6 vehicles (V1--V6). The true box of V2 has corners $(100,100)$--$(140,130)$ pixels. The detector returns 9 boxes, sorted by confidence, each listed with the vehicle it overlaps most and that IoU:

\begin{center}\small\begin{tabular}{@{}cccccccccccc@{}}\toprule Rank&1&2&3&4&5&6&7&8&9\\
Confidence&0.97&0.93&0.90&0.85&0.80&0.74&0.60&0.52&0.41\\
Best match&V1&V2&V1&V3&V4&V5&none&V4&V6\\
IoU&0.81&?&0.55&0.72&0.30&0.58&0&0.68&0.45\\\bottomrule
\end{tabular}\end{center}

(a) Detection 2 has corners $(105,98)$--$(147,128)$. Compute its IoU with V2. [4]\newline (b) At IoU threshold 0.5, label each detection TP or FP, remembering that each true vehicle can be matched only once. [4]\newline (c) Compute running precision and recall, and the interpolated average precision. [8]\newline (d) Recompute AP at a stricter IoU threshold of 0.7 and explain what the drop tells an imagery analyst. [4]\end{prob}

\begin{soln}\textbf{(a)} Overlap in $x$: $105$ to $140$ = 35 px; in $y$: $100$ to $128$ = 28 px; intersection $=980$. Areas: true $40\times30=1{,}200$, predicted $42\times30=1{,}260$. Union $=1{,}200+1{,}260-980=1{,}480$. IoU $=980/1{,}480=\mathbf{0.66}$.

\textbf{(b)} TP, TP, \textbf{FP} (V1 already matched by detection 1: a duplicate), TP, \textbf{FP} (IoU 0.30), TP, \textbf{FP} (no vehicle: decoy or clutter), TP (V4 at 0.68 --- V4 was not matched by detection 5), \textbf{FP} (IoU 0.45 $<$ 0.5). Five TPs; V6 is never found.

\textbf{(c)} Precision: 1.00, 1.00, 0.67, 0.75, 0.60, 0.67, 0.57, 0.63, 0.56. Recall (of 6): 0.17, 0.33, 0.33, 0.50, 0.50, 0.67, 0.67, 0.83, 0.83. Interpolated precision at each recall step (the best precision at that recall or beyond): 1.00 at 0.17 and 0.33; 0.75 at 0.50; 0.67 at 0.67; 0.63 at 0.83. AP $=\frac16(1.00+1.00+0.75+0.67+0.63)=\mathbf{0.67}$. Recall never reaches 1, so the last sixth of the curve contributes nothing.

\textbf{(d)} At 0.7 only detections 1 (0.81) and 4 (0.72) count. Precision 1.00, 0.50, 0.33, 0.50, 0.40, \ldots{}; recall reaches only $2/6$. AP $=\frac16(1.00)+\frac16(0.50)=\mathbf{0.25}$. The detector usually \textit{finds} vehicles but localises them loosely. For counting and cueing that may be acceptable; for precise targeting or measuring vehicle dimensions it is not. Report AP at the IoU the mission needs, and note that the decoy (detection 7) scored 0.60 confidence --- a deception risk to test explicitly.\end{soln}

\end{document}